\scanpage{24}
On vérifie que les $\pi^n:\Gamma^+(M)\longrightarrow\Gamma^+M$ ainsi construits satisfont les conditions 1.1. à 1.6. ($\Gamma^+(M)$ considéré comme un idéal de $\Gamma(M)$, ou comme une $\Lambda$-algèbre).

$\Gamma^{\unsure{+|*}}(M)$ devient ainsi une $\Lambda$-algèbre à puissances divisées. Elle est munie d'un hom de $\Lambda$-modules
\[
M\longrightarrow\Gamma^+(M)\ ;
\]
ce \unsure{dernier} \unsure{joue} un rôle \uline{universel}.

\uline{§ 3} Soit $A$ un anneau, et $I$ un idéal de $A$. On s'intéresse, --- pour un anneau $B$, muni d'un idéal $J$ \uline{à puissances divisées}, aux homs d'anneaux
\[
u:A\longrightarrow B
\]
tels que $u(I)\subset J$. On cherche un tel \unsure{fourbi|fonct.}, \uline{universel}. On prend l'\deleted{\unsure{anneau}} \inserted{$A$-algèbre} $\Gamma^{\unsure{+|*}}(A,\unsure{I|J})$ engendré par des éléments
\[
\pi^n x\qquad \deleted{n\in\unsure{\mathbb N^{*}|\mathbb N}}\inserted{n\geq2},\ x\in\unsure{I|J}
\]
avec les relations 1.2. à 1.6. (où \deleted{\unsure{i}} $x,y\in\unsure{I|J}$, $\lambda\in A$, et $\pi^1(x)$ doit être $x$). Une autre façon de le dire est celle-ci: $\unsure{I|J}$ est un $A$-module, on construit $\Gamma_A^+(\unsure{J|I})$, et on
\mspage{23}
